\appendix

\chapter{Tablas Detalladas de Resultados}

\section{Hiperparámetros Óptimos por Escenario}

La Tabla~\ref{tab:hyperparams_trss} presenta los hiperparámetros óptimos encontrados por Optuna para TRSS en cada combinación de dataset y escenario de degradación. Estos valores son reproducibles ejecutando \texttt{experiments/run\_optuna\_optimization.py}.

\begin{table}[ht]
\centering
\caption{Hiperparámetros óptimos de TRSS por dataset y escenario representativo.}
\label{tab:hyperparams_trss}
\begin{tabular}{llcccc}
\toprule
\textbf{Dataset} & \textbf{Escenario} & $k$ & $\sigma$ & $\alpha$ & $\beta$ \\
\midrule
BCI IV 2a  & nearby, 1 canal  & 6  & 0.856 & 0.0114 & 0.765 \\
BCI IV 2a  & nearby, 10\%     & 5  & 0.712 & 0.0230 & 0.612 \\
BCI IV 2a  & nearby, 40\%     & 4  & 0.950 & 0.0085 & 0.890 \\
MNE Sample & nearby, 10\%     & 8  & 1.200 & 0.0150 & 0.550 \\
MNE Sample & nearby, 40\%     & 6  & 0.900 & 0.0200 & 0.780 \\
PhysioNet  & random, 10\%     & 10 & 1.500 & 0.0300 & 0.420 \\
PhysioNet  & nearby, 40\%     & 7  & 1.100 & 0.0120 & 0.850 \\
\bottomrule
\end{tabular}
\end{table}

Patrones observados en la sintonización:
\begin{itemize}
    \item $\beta$ (peso temporal) tiende a incrementarse con la severidad de la pérdida, confirmando que la regularización temporal se vuelve más valiosa cuanta más información espacial se pierde.
    \item $\alpha$ (peso espacial) se mantiene bajo ($< 0.05$), indicando que un suavizado espacial excesivo es contraproducente.
    \item $k$ (vecinos del grafo) se adapta a la densidad del dataset: valores más altos para PhysioNet (64 ch) y más bajos para BCI IV (22 ch).
\end{itemize}

\section{Métricas Espectrales Detalladas}

\begin{table}[ht]
\centering
\caption{LSD y Coherencia promedio por método y dataset (valores de \texttt{optuna\_summary\_pivot.csv}).}
\label{tab:spectral_detail}
\begin{tabular}{lcccccc}
\toprule
& \multicolumn{3}{c}{\textbf{LSD} $\downarrow$} & \multicolumn{3}{c}{\textbf{Coherencia} $\uparrow$} \\
\cmidrule(lr){2-4} \cmidrule(lr){5-7}
\textbf{Método} & \textbf{BCI IV} & \textbf{MNE} & \textbf{PhysioNet} & \textbf{BCI IV} & \textbf{MNE} & \textbf{PhysioNet} \\
\midrule
TRSS         & 0.124 & 0.111 & 0.189 & 0.997 & 0.994 & 0.990 \\
Tikhonov     & 0.314 & 0.460 & 0.308 & 0.965 & 0.916 & 0.973 \\
RBF (TPS)    & 0.346 & 0.662 & 0.328 & 0.955 & 0.835 & 0.968 \\
Sph. Spline  & 0.369 & 0.695 & 0.382 & 0.950 & 0.816 & 0.946 \\
ICA          & 0.591 & 0.465 & 0.621 & 0.902 & 0.917 & 0.896 \\
\bottomrule
\end{tabular}
\end{table}

\chapter{Estructura del Repositorio}

El código fuente y los artefactos experimentales se organizan bajo el directorio \texttt{Thesis-Copilot-Toolkit/} de la siguiente manera:

\begin{verbatim}
Thesis-Copilot-Toolkit/
+-- src/
|   +-- data/
|   |   +-- data_loader.py          # Cargadores de datasets
|   +-- graph_construction/
|   |   +-- graph_constructors.py   # 10 constructores de grafos
|   |   +-- aew.py                  # Adaptive Edge Weighting
|   +-- interpolation_methods.py    # 18 algoritmos
|   +-- evaluation.py              # 6 metricas
+-- experiments/
|   +-- run_optuna_optimization.py  # Optimizacion GSP/temporal
|   +-- run_optuna_baselines.py     # Optimizacion baselines
|   +-- plot_*.py                   # Scripts de visualizacion
|   +-- finalize_b3_b4.py          # Consolidacion final
+-- results_optuna_final/
|   +-- optuna_best_results.csv     # Resultados optimos GSP
|   +-- optuna_summary_pivot.csv    # Tabla pivote consolidada
|   +-- degradation_*.png           # Curvas de degradacion
|   +-- heatmap_*.png               # Mapas de calor
|   +-- radar_*.png                 # Graficos radar
|   +-- erp_timeseries_*.png        # Series ERP
|   +-- psd_comparison_*.png        # Comparaciones PSD
+-- results_optuna_baselines/
|   +-- optuna_best_results_baselines.csv
|   +-- erp_optimized_baselines_*.png
|   +-- psd_optimized_baselines_*.png
+-- shared/
|   +-- results_data.tex            # Variables LaTeX compartidas
+-- thesis/usm/                     # Documento de tesis
+-- paper/ieee/                     # Documento del paper
\end{verbatim}

\chapter{Instrucciones de Reproducción}

Para reproducir los experimentos reportados en esta tesis, ejecutar los siguientes comandos desde el directorio raíz del repositorio:

\begin{enumerate}
    \item \textbf{Instalar dependencias:}
\begin{verbatim}
pip install -r requirements.txt
\end{verbatim}

    \item \textbf{Configurar rutas de datasets:}
\begin{verbatim}
export BCI_IV_2A_PATH=./datasets/BCICIV_2a_gdf
export EEGBCI_LOCAL_PATH=./datasets/MNE-eegbci-data
\end{verbatim}

    \item \textbf{Ejecutar optimización de métodos GSP/temporales:}
\begin{verbatim}
python experiments/run_optuna_optimization.py
\end{verbatim}

    \item \textbf{Ejecutar optimización de baselines:}
\begin{verbatim}
python experiments/run_optuna_baselines.py
\end{verbatim}

    \item \textbf{Generar figuras:}
\begin{verbatim}
python experiments/plot_degradation_curves.py
python experiments/plot_radar_charts.py
python experiments/plot_heatmaps.py
python experiments/plot_erp_time_series.py
python experiments/plot_optimized_baselines.py
\end{verbatim}

    \item \textbf{Compilar documentos:}
\begin{verbatim}
cd Thesis-Copilot-Toolkit/thesis/usm && latexmk -pdf main.tex
cd Thesis-Copilot-Toolkit/paper/ieee && latexmk -pdf main.tex
\end{verbatim}
\end{enumerate}

Todas las operaciones estocásticas utilizan semillas fijas (\texttt{random\_state=42}) para garantizar reproducibilidad determinista.
